\textbf{1. Kiến thức cần nhớ}
\begin{itemize}
\item \textbf{Phép chia đơn thức cho đơn thức}:
\begin{itemize}
\item Đơn thức $ chia hết cho đơn thức $ (với  \ne 0$) khi mỗi biến của $ đều là biến của $ với số mũ không lớn hơn số mũ của nó trong $.
\item Quy tắc chia: Ta chia hệ số của $ cho hệ số của $, chia lũy thừa của từng biến trong $ cho lũy thừa của cùng biến đó trong $, rồi nhân các kết quả lại:
80493\boxed{x^m : x^n = x^{m-n}\quad (m \ge n)}80493
\end{itemize}
\item \textbf{Phép chia đa thức cho đơn thức}:
\begin{itemize}
\item Đa thức $ chia hết cho đơn thức $ nếu mọi hạng tử của $ đều chia hết cho $.
\item Quy tắc chia: Muốn chia đa thức $ cho đơn thức $, ta chia từng hạng tử của $ cho $ rồi cộng các kết quả lại:
80493\boxed{(A + B - C) : D = A : D + B : D - C : D}80493
\end{itemize}
\end{itemize}

\textbf{2. Phương pháp giải các dạng}

%% DANG 8D1?5-1
\textbf{Dạng 1. Chia đơn thức cho đơn thức}
\begin{enumerate}
\item Chia phần hệ số: lấy hệ số đơn thức bị chia chia cho hệ số đơn thức chia.
\item Chia phần biến: với mỗi biến chung, lấy số mũ của số bị chia trừ số mũ của số chia (^m : x^n = x^{m-n}$); quy ước ^0 = 1$.
\item Nhân các kết quả thành phần lại với nhau.
\end{enumerate}
Lưu ý: Biến nào có số mũ bằng nhau thì kết quả chia bằng $ (ví dụ: ^2 : y^2 = 1$).

%% DANG 8D1?5-2
\textbf{Dạng 2. Chia đa thức cho đơn thức}
\begin{enumerate}
\item Tách phép chia đa thức thành các phép chia từng đơn thức: lấy mỗi hạng tử của đa thức chia cho đơn thức.
\item Thực hiện phép chia đơn thức cho đơn thức đối với từng cặp (chú ý giữ đúng dấu của mỗi hạng tử).
\item Cộng các đơn thức kết quả lại và thu gọn (nếu có đồng dạng).
\end{enumerate}

%% DANG 8D1?5-3
\textbf{Dạng 3. Xét tính chia hết, tìm điều kiện để phép chia hết}
\begin{enumerate}
\item \textbf{Xét tính chia hết}:
\begin{itemize}
\item Với đơn thức $ chia cho đơn thức $: kiểm tra xem mọi biến của $ có mặt trong $ với số mũ nhỏ hơn hoặc bằng số mũ trong $ hay không.
\item Với đa thức chia cho đơn thức: kiểm tra xem tất cả các hạng tử của đa thức có cùng chia hết cho đơn thức chia hay không.
\end{itemize}
\item \textbf{Tìm điều kiện của số tự nhiên $}:
\begin{itemize}
\item Thiết lập các bất đẳng thức về số mũ của từng biến: số mũ trong đơn thức bị chia $\ge$ số mũ trong đơn thức chia.
\item Giải các bất đẳng thức tìm giá trị số tự nhiên $ thỏa mãn đồng thời tất cả các điều kiện.
\end{itemize}
\end{enumerate}
